\section{Partie théorique} \label{sec:theorie}

\subsection{Oscillations libres d’un pendule de Pohl} \label{subsec:osc_libres}

Le pendule de Pohl est un oscillateur mécanique rotatif composé d'un disque couplé à un ressort
spiral. La coordonnée généralisée du système est l’angle de torsion $\theta(t)$.
En appliquant le théorème du moment cinétique au disque, on obtient :

\begin{equation}
\frac{dL}{dt} = \sum M_{\mathrm{ext}},
\qquad L = J\,\dot{\theta},
\label{eq:TMC}
\end{equation}

où $J$ est le moment d’inertie du balancier. Deux moments mécaniques agissent sur le système :
\begin{equation}
M_r = -C\,\theta,
\qquad
M_f = -b\,\dot{\theta},
\label{eq:MOMENTS}
\end{equation}
où $C$ est la constante de torsion, et $b$ le moment de frottement visqueux.  
En combinant ces contributions, l'équation du mouvement devient :

\begin{equation}
J\ddot{\theta} + b\dot{\theta} + C\theta = 0.
\label{eq:edo}
\end{equation}

Dont on définit la pulsation propre non amortie $\omega_0$ et le coefficient d'amortissement $\lambda$ tels que :
\begin{equation}
\omega_0 = \sqrt{\frac{C}{J}},
\qquad
\lambda = \frac{b}{2J}.
\label{eq:lambda} % et omega 0
\end{equation}

L’équation différentielle se met sous la forme caractéristique :
\begin{equation}
\ddot{\theta} + 2\lambda \dot{\theta} + \omega_0^2 \theta = 0.
\label{eq:cara}
\end{equation}

Dans le régime faiblement amorti ($\lambda < \omega_0$), la solution générale est :

\begin{equation}
\theta(t) = \theta_0\,e^{-\lambda t}
\cos(\omega_1 t + \phi),
\label{eq:solution}
\end{equation}

où la pseudo-pulsation est :

\begin{equation}
\omega_1 = \sqrt{\omega_0^2 - \lambda^2}.
\label{eq:w}
\end{equation}

La période mesurée expérimentalement est alors :

\begin{equation}
T = \frac{2\pi}{\omega_1}.
\label{eq:t}
\end{equation}

L’expression de $J$ peut être inversée en fonction des paramètres mesurés :

\begin{equation}
J = \frac{C}{\omega_0^2}.
\label{eq:J}
\end{equation}


\subsection{Oscillations forcées} \label{subsec:osc_forcees}

Lorsque le moteur excite le système, un moment sinusoïdal est appliqué :

\begin{equation}
M_e(t) = M_0 \cos(\omega t),
\label{eq:moment_excitation}
\end{equation}

ce qui modifie l'équation différentielle :

\begin{equation}
J\ddot{\theta} + b\dot{\theta} + C\theta = M_0\cos(\omega t).
\label{eq:equa_forcee_brut}
\end{equation}

En divisant par $J$, on obtient l'équation caractéristique sous forme normalisée :

\begin{equation}
\ddot{\theta} + 2\lambda \dot{\theta} + \omega_0^2 \theta
= A\cos(\omega t),
\qquad A = \frac{M_0}{J}.
\label{eq:equa_forcee_normalisee}
\end{equation}


\subsection{Méthode des phaseurs} \label{subsec:phaseurs}

On cherche une solution sinusoïdale particulière de la forme :

\begin{equation}
\theta(t) = \Theta(\omega)\cos(\omega t + \phi(\omega)).
\label{eq:solution_phaseurs}
\end{equation}

En utilisant la méthode complexe (phaseurs), l'amplitude en régime forcé est :

\begin{equation}
\Theta(\omega)
= \frac{A}{\sqrt{(\omega_0^2 - \omega^2)^2 + (2\lambda\omega)^2}}.
\label{eq:amplitude_forcee}
\end{equation}

La phase du mouvement par rapport à la force excitatrice est :

\begin{equation}
\phi(\omega)
= \arctan\!\left(\frac{2\lambda\omega}{\omega_0^2 - \omega^2}\right).
\label{eq:phase_forcee}
\end{equation}

La résonance apparaît lorsque l’amplitude est maximale, ce qui se produit pour :

\begin{equation}
\omega_{\mathrm{res}} = \sqrt{\omega_0^2 - 2\lambda^2}.
\label{eq:omega_res}
\end{equation}


\subsection{Facteur de qualité} \label{subsec:facteur_Q}

Le facteur de qualité d’un oscillateur amorti est défini par :

\begin{equation}
Q = \frac{\omega_0}{2\lambda}.
\label{eq:Q_lambda}
\end{equation}

Dans l’analyse fréquentielle (courbe de puissance), il peut aussi être obtenu via :

\begin{equation}
Q = \frac{f_0}{\mathrm{FWHM}},
\label{eq:Q_fwhm}
\end{equation}

où $\mathrm{FWHM}$ est la largeur à mi-hauteur du pic de résonance,
et $f_0 = \omega_0 / (2\pi)$ la fréquence propre.

